
\clearpage
\part{A refresher on computability, from the paper's viewpoint}
\label{part:refresher}
\setlength{\parskip}{0.30em}

The outermost layer of the argument is classical computation.  Before any
low-degree test or entangled strategy appears, the proof must be able to take
the finite text of one verifier, transform that text into another verifier,
and arrange for a verifier to use a description of itself.  This part rebuilds
that vocabulary from the ground up.  It is deliberately qualitative about
polynomial overheads; the exact costed statements begin in Part~II.
Figure~\ref{fig:pipeline-glance} gives the whole description-level pipeline
and its fixed-point feedback at a glance, while Figure~\ref{fig:roadmap} maps
the refresher objects to their resource-accounted counterparts.

\section{Turing machines in the form used by the paper}
\label{sec:refresher-tm}

\begin{figure}[h]
\centering
\begin{tikzpicture}[node distance=2.6mm and 8mm]
  \node[ctrl, minimum width=35mm, minimum height=16mm] (control)
    {finite control $q\in Q$\\[1pt]
     $\delta:Q\!\times\!\Gamma^{k+2}\to Q\!\times\!\Gamma^2\!\times\!\{L,S,R\}^{k+2}$};
  \node[lbl, above=2mm of control, text=machine]
    {$M=(Q,\Sigma,\Gamma,\delta,q_0,q_{\mathrm h})$};

  \node[tapelbl, below left=10mm and 8mm of control.south] (l1) {input $x_1$};
  \node[cell head, right=1mm of l1] (i11) {$1$};
  \node[cell, right=0pt of i11] (i12) {$0$};
  \node[cell blank, right=0pt of i12] (i13) {$\sqcup$};
  \node[cell blank, right=0pt of i13] (i14) {$\sqcup$};
  \node[head, below=.8mm of i11] {};

  \node[tapelbl, below=of l1] (l2) {input $x_2$};
  \node[cell, right=1mm of l2] (i21) {$1$};
  \node[cell head, right=0pt of i21] (i22) {$0$};
  \node[cell blank, right=0pt of i22] (i23) {$\sqcup$};
  \node[cell blank, right=0pt of i23] (i24) {$\sqcup$};
  \node[head, below=.8mm of i22] {};

  \node[lbl, below=1mm of l2] (dots) {$\vdots$};
  \node[tapelbl, below=1mm of dots] (lk) {input $x_k$};
  \node[cell head, right=1mm of lk] (ik1) {$1$};
  \node[cell, right=0pt of ik1] (ik2) {$1$};
  \node[cell, right=0pt of ik2] (ik3) {$0$};
  \node[cell blank, right=0pt of ik3] (ik4) {$\sqcup$};
  \node[head, below=.8mm of ik1] {};

  \node[tapelbl, below=7mm of lk] (lw) {work};
  \node[cell, right=1mm of lw] (w1) {$1$};
  \node[cell head, right=0pt of w1] (w2) {$0$};
  \node[cell blank, right=0pt of w2] (w3) {$\sqcup$};
  \node[cell blank, right=0pt of w3] (w4) {$\sqcup$};
  \node[head, below=.8mm of w2] {};

  \node[tapelbl, below=of lw] (lo) {output};
  \node[cell head, right=1mm of lo] (o1) {$1$};
  \node[cell blank, right=0pt of o1] (o2) {$\sqcup$};
  \node[cell blank, right=0pt of o2] (o3) {$\sqcup$};
  \node[cell blank, right=0pt of o3] (o4) {$\sqcup$};
  \node[head, below=.8mm of o1] {};

  \node[checked, right=22mm of i22, text width=29mm] (read)
    {read-only\\$T_1,\ldots,T_k$};
  \node[checked, right=22mm of w2, text width=29mm] (write)
    {writable\\$T_{k+1},T_{k+2}$};
  \draw[->, machine, line width=.8pt] (i14.east) -| (control.south east);
  \draw[->, machine, line width=.8pt] (ik4.east) -| (control.south east);
  \draw[->, machine, line width=.8pt] (control.south west) |- (w4.east);
  \draw[->, machine, line width=.8pt] (control.south west) |- (o4.east);
  \draw[callout] (read) -- (i24.east);
  \draw[callout] (write) -- (w4.east);
  \node[lblsm, below=3mm of o1.south west, anchor=west]
    {amber cell $=$ symbol scanned now \qquad triangle $=$ head position};
\end{tikzpicture}

\caption{Amber cells are the symbols read by the finite control, with arrows returning updates only to the work and output tapes.  This concrete separation of finite rule book from finitely used memory is the machine model simulated throughout the document.}
\label{fig:tm-anatomy}
\end{figure}


\subsection{Programs, configurations, and one step}

\begin{figure}[t]
\centering
\begin{tikzpicture}[node distance=3mm and 5mm]
  % ---- BEFORE: the state sits beside its own tape, on one baseline, so that
  % nothing has to be pointed at across the picture.
  \node[state initial] (q) {$q$};
  \node[lbl, text=machine, font=\footnotesize\sffamily\bfseries,
        above=2mm of q.north west, anchor=south west] {BEFORE};
  \node[cell blank, right=4mm of q] (b0) {$\sqcup$};
  \node[cell head, right=0pt of b0] (b1) {$0$};
  \node[cell, right=0pt of b1] (b2) {$1$};
  \node[cell blank, right=0pt of b2] (b3) {$\sqcup$};
  \node[cell blank, right=0pt of b3] (b4) {$\sqcup$};
  \node[head, below=.8mm of b1] (hb) {};
  \node[lblsm, right=0.5mm of hb] {$h$};

  % ---- the rule that is applied once
  \node[redex, right=9mm of b4, minimum size=13mm] (delta) {$\delta$};
  \node[lbl, above=3mm of delta] {$\delta(q,0)=(q',1,R)$};
  \draw[xform] (b4.east) -- (delta.west);

  % ---- AFTER
  \node[state, fill=focuslt, draw=focus, right=9mm of delta] (qp) {$q'$};
  \node[lbl, text=machine, font=\footnotesize\sffamily\bfseries,
        above=2mm of qp.north west, anchor=south west] {AFTER};
  \draw[xform] (delta.east) -- (qp.west);
  \node[cell blank, right=4mm of qp] (a0) {$\sqcup$};
  \node[cell, fill=focuslt, draw=focus, line width=.9pt, right=0pt of a0] (a1) {$1$};
  \node[cell head, right=0pt of a1] (a2) {$1$};
  \node[cell blank, right=0pt of a2] (a3) {$\sqcup$};
  \node[cell blank, right=0pt of a3] (a4) {$\sqcup$};
  \node[head, below=.8mm of a2] (ha) {};
  % beside, not below, the head marker: below it the write callout would
  % have to be crossed.
  \node[lblsm, right=0.5mm of ha] (hpl) {$h+1$};

  % ---- the three local changes, listed in the same left-to-right order as
  % their targets so that no two callouts cross.
  \node[checked, below=13mm of qp.south, minimum width=22mm] (statechange)
    {state: $q\mapsto q'$};
  \node[checked, right=3mm of statechange, minimum width=22mm] (writechange)
    {write: $0\mapsto1$};
  \node[checked, right=3mm of writechange, minimum width=26mm] (headchange)
    {head: $h\mapsto h+1$};
  \draw[callout] (statechange.north) -- (qp.south);
  \draw[callout] (writechange.north) -- (a1.south);
  \draw[callout] (headchange.north) -- (ha.south);
  \node[lblsm, below=2mm of writechange]
    {every other tape cell is unchanged};
\end{tikzpicture}

\caption{The amber transition rewrites the scanned zero, changes \(q\) to \(q'\), and advances the head from \(h\) to \(h+1\).  Exposing exactly these three local changes is what later permits a computation history to be checked by bounded windows.}
\label{fig:one-step}
\end{figure}


A Turing machine separates two things which ordinary mathematical notation
often blends together: a finite rule book and an unbounded but finitely used
memory.  For one tape, a convenient formal tuple is
\[
 M=(Q,\Sigma,\Gamma,\delta,q_0,q_{\mathrm h}),
\]
where \(Q\) is a finite set of control states, \(q_0\in Q\) is the initial
state, \(q_{\mathrm h}\in Q\) is the halting state,
\(\Sigma=\{0,1\}\) is the input alphabet,
\(\Gamma\supseteq\Sigma\cup\{\sqcup\}\) is a finite tape alphabet with blank
symbol \(\sqcup\notin\Sigma\), and
\[
 \delta:(Q\setminus\{q_{\mathrm h}\})\times\Gamma
 \longrightarrow Q\times\Gamma\times\{L,S,R\}
\]
is the transition function.  We take tapes to be indexed by \(\N\), so a
left move at cell zero leaves the head at zero.  The paper uses a multitape
version: a \(k\)-input machine has \(k\) input tapes, one work tape, and one
output tape.  Its transition reads all scanned symbols, changes the state,
writes on the work and output tapes, and moves each head by at most one cell.
The input tapes are read-only.  Adding separate accept and reject halting
states, or allowing writes on input tapes, changes none of the computability
or polynomial-overhead statements below.
Figure~\ref{fig:tm-anatomy} fixes this multitape anatomy and the semantic roles
of its heads, tapes, and finite control.

A \emph{configuration} is a complete instantaneous description
\[
 C=(q,T_1,h_1,\ldots,T_{k+2},h_{k+2}):
\]
the current state, the contents \(T_i:\N\to\Gamma\) of every tape, and the
head locations \(h_i\in\N\).  Only finitely many cells are nonblank in every
reachable configuration.  The notation \(C\vdash_M C'\) means that one
application of \(\delta\) changes \(C\) into \(C'\).  Thus a computation is
not a metaphor: it is the uniquely determined path
\(C_0\vdash_M C_1\vdash_M C_2\vdash_M\cdots\) in the graph of configurations.

\paragraph{Why this matters here.}
Later, Cook--Levin replaces a bounded computation by local constraints on
adjacent configurations.  That replacement works because one step looks only
at the finite control and the cells under the heads.  The enormous object is
the history; the local dynamical law \(\delta\) is finite.  This is close to a
lattice model with a fixed local update rule, except that the update is
deterministic and the head positions are part of the state.

\paragraph{Worked one-step example.}
If \(\delta(q,0)=(q',1,R)\), then
\[
 (q,\ldots\sqcup\,\underline 0\,1\sqcup\ldots,h)
 \ \vdash_M\ 
 (q',\ldots\sqcup\,1\,\underline 1\sqcup\ldots,h+1).
\]
Only the state, the scanned cell, and the head location have changed.
Figure~\ref{fig:one-step} isolates those three changes in a literal tape move.

The machine \emph{halts} on input \(x\) if some \(C_T\) has state
\(q_{\mathrm h}\).  Its output is the longest nonblank initial string on the
output tape.  We write \(M(x)\) for that string and \(M(x)=\bot\) if the path
never reaches \(q_{\mathrm h}\).  The instance running time is
\[
 \Time_{M,x}:=\min\{T:C_T\text{ has state }q_{\mathrm h}\},
\]
with value \(\infty\) if the set is empty.  We sometimes abbreviate this as
\(\Time_M(x)\).  Notice the distinction between \emph{has not halted yet} and
\emph{will never halt}; no finite observation of a run decides the latter in
general.
Figure~\ref{fig:config-path} contrasts the finite witness supplied by a
halting configuration with an open-ended observed prefix.

\begin{figure}[t]
\centering
\begin{tikzpicture}[node distance=5mm and 7mm]
  \node[lbl, text=machine, font=\footnotesize\sffamily\bfseries] (haltlabel)
    {A HALTING RUN};
  \node[config, below=3mm of haltlabel] (h0) {$C_0$};
  \node[config, right=of h0] (h1) {$C_1$};
  \node[config, right=of h1] (h2) {$C_2$};
  \node[lbl, right=of h2] (hdots) {$\cdots$};
  \node[config, right=of hdots] (htm) {$C_{T-1}$};
  \node[state halt, right=of htm, fill=focuslt, draw=focus] (ht) {$q_{\mathrm h}$};
  \foreach \a/\b in {h0/h1,h1/h2,h2/hdots,hdots/htm,htm/ht}
    \draw[->, machine, line width=.8pt] (\a) -- (\b);
  \node[checked, below=4mm of ht, minimum width=30mm] (time)
    {$\Time_{M,x}=T$};
  \draw[callout] (time) -- (ht);

  \node[lbl, text=machine, font=\footnotesize\sffamily\bfseries,
        below=16mm of h0.south, anchor=west] (openlabel)
    {AN OBSERVED PREFIX};
  \node[config, below=3mm of openlabel] (o0) {$C_0$};
  \node[config, right=of o0] (o1) {$C_1$};
  \node[config, right=of o1] (o2) {$C_2$};
  \node[lbl, right=of o2] (odots) {$\cdots$};
  \node[config, right=of odots, fill=focuslt, draw=focus] (ot) {$C_t$};
  \node[config, right=of ot] (otp) {$C_{t+1}$};
  \node[lbl, right=of otp, text=machine] (horizon) {$\cdots\ \longrightarrow$};
  \foreach \a/\b in {o0/o1,o1/o2,o2/odots,odots/ot,ot/otp,otp/horizon}
    \draw[->, machine, line width=.8pt] (\a) -- (\b);
  \draw[machine, line width=.7pt, opacity=.55] (horizon.east) -- ++(12mm,0);
  \draw[brace] ([yshift=-4mm]o0.south west) -- ([yshift=-4mm]ot.south east)
    node[midway, below=4mm, lblsm] {checked through stage $t$};
  \node[bad, rounded corners=2pt, right=9mm of ot, yshift=-13mm, inner sep=4pt,
        font=\scriptsize\sffamily] (never) {not a proof of \emph{never halts}};
  \draw[callout] (never.north) -- (ot.south);
\end{tikzpicture}

\caption{The upper deterministic path ends in \(q_{\mathrm h}\) and certifies the exact time \(T\), whereas the lower path fades beyond the last observed configuration \(C_t\).  The missing halt in a finite prefix cannot be promoted to a proof that the run never halts.}
\label{fig:config-path}
\end{figure}

\paragraph{Why this matters here.}
The compression argument repeatedly says that a generated verifier finishes
within a supplied bound.  That is a statement about the number of transitions
on every relevant input, not merely about the mathematical function obtained
after ignoring time.  A timeout wrapper may safely reject after \(T\) steps,
because it is counting the very transitions used in the definition above.

\paragraph{Worked timing example.}
For the transition just displayed, if the new state \(q'\) is the halt state,
then the initial configuration is assigned time one.  A timeout of zero
rejects without applying the transition; a timeout of one obtains the output.

\subsection{A complete two-state equality machine}

Here is a small instance in the paper's multitape style.  The machine has two
one-bit input tapes \(A,B\), a work tape \(W\), and an output tape \(O\).  It
has two nonhalting states \(q_0,q_1\), plus the designated state
\(q_{\mathrm h}\).  A table entry lists
\((\hbox{next state};\hbox{new }W,\hbox{new }O;
  \hbox{moves of }A,B,W,O)\).
Rows are tested from top to bottom; they are disjoint except for the last
catch-all row, so the table specifies every possible scanned tuple.

\begin{center}
\small
\begin{tabularx}{\textwidth}{P{0.09\textwidth}P{0.24\textwidth}X}
\toprule
State & Scanned \((A,B,W,O)\) & Transition \\
\midrule
\(q_0\) & \((0,0,\sqcup,\sqcup)\) or
          \((1,1,\sqcup,\sqcup)\)
 & \((q_1;1,\sqcup;R,R,S,S)\) \\
\(q_0\) & \((0,1,\sqcup,\sqcup)\) or
          \((1,0,\sqcup,\sqcup)\)
 & \((q_1;0,\sqcup;R,R,S,S)\) \\
\(q_1\) & \((\sqcup,\sqcup,r,\sqcup)\), \(r\in\{0,1\}\)
 & \((q_0;r,r;S,S,S,S)\) \\
\(q_0\) & \((\sqcup,\sqcup,r,r)\), \(r\in\{0,1\}\)
 & \((q_{\mathrm h};r,r;S,S,S,S)\) \\
\(q_0,q_1\) & every tuple not covered above
 & \((q_{\mathrm h};0,0;S,S,S,S)\) \\
\bottomrule
\end{tabularx}
\end{center}

Figure~\ref{fig:eq-machine-states} turns the five rows into a state diagram,
including the totalizing catch-all path.

\begin{figure}[t]
\centering
% local: every edge label carries the page fill so that it interrupts, rather
% than is crossed by, the edge it names.  The two long q0 -> qh edges are
% routed orthogonally above and below the q0/q1 traffic instead of bowing
% through the labels that sit there.
\begin{tikzpicture}[node distance=18mm and 38mm,
  elbl/.style={lblsm, fill=paper, inner sep=2pt, align=center},
  rulenum/.style={text=focus, font=\scriptsize\sffamily\bfseries}]
  \node[state initial] (q0) {$q_0$};
  \node[state, right=56mm of q0] (q1) {$q_1$};
  \node[state halt, right=52mm of q1, minimum size=9.5mm] (qh) {$q_{\mathrm h}$};
  \draw[->, machine, line width=.8pt] ([xshift=-11mm]q0.west) -- (q0)
    node[lblsm, above, pos=.05] {start};

  % 1, 2: the two legal first steps; 3: the copy step back to q0
  \draw[->, machine, line width=.8pt]
    (q0) to[bend left=32] node[elbl, pos=.5]
    {\textcolor{focus}{\bfseries 1}\; equal inputs; $W\leftarrow1$; $RRSS$} (q1);
  \draw[->, machine, line width=.8pt]
    (q0) -- node[elbl, pos=.5, below=0pt]
    {\textcolor{focus}{\bfseries 2}\; unequal inputs; $W\leftarrow0$; $RRSS$} (q1);
  \draw[->, machine, line width=.8pt]
    (q1) to[bend left=32] node[elbl, pos=.5]
    {\textcolor{focus}{\bfseries 3}\; copy; $O\leftarrow r$; $SSSS$} (q0);

  % 4: the halt row, routed over the q0/q1 traffic
  \draw[->, machine, line width=.8pt, rounded corners=3pt]
    (q0.north) -- ++(0,21mm) -| node[elbl, pos=.35]
    {\textcolor{focus}{\bfseries 4}\; $W,O$ agree; halt; $SSSS$} (qh.north);
  % 5: the catch-all row, which fires from either nonhalting state
  \draw[->, machine, line width=.8pt, rounded corners=3pt]
    (q0.south) -- ++(0,-21mm) -| node[elbl, pos=.35]
    {\textcolor{focus}{\bfseries 5}\; any other tuple; $W,O\leftarrow0$; $SSSS$}
    (qh.south);
  \draw[->, machine, line width=.8pt]
    (q1) -- node[elbl, pos=.5] {\textcolor{focus}{\bfseries 5}} (qh);
\end{tikzpicture}

\caption{The five numbered paths are exactly the equal-input, unequal-input, copy, halt, and catch-all rows of \(M_=\).  Reading the table as control flow makes totality and the three-step legal-input route visible at once.}
\label{fig:eq-machine-states}
\end{figure}

The final row makes the table total and rejects malformed inputs.  On legal
one-bit inputs the first step computes the equality bit on \(W\), the second
copies it to \(O\), and the third enters the halt state.  For \(A=B=1\), the
full three-step trace, showing the scanned cells, is
\[
\begin{array}{c|c|c|c|c|c|l}
i&q&A&B&W&O&\text{rule used}\\ \hline
0&q_0&\underline1&\underline1&\underline\sqcup&\underline\sqcup
  &\text{equal-input row}\\
1&q_1&1\underline\sqcup&1\underline\sqcup&\underline1&\underline\sqcup
  &\text{copy row with }r=1\\
2&q_0&1\underline\sqcup&1\underline\sqcup&\underline1&\underline1
  &\text{halt row with }r=1\\
3&q_{\mathrm h}&1\underline\sqcup&1\underline\sqcup&\underline1&\underline1
  &\text{halted.}
\end{array}
\]
Thus \(M_{=}(1,1)=1\) and \(\Time_{M_=,(1,1)}=3\).  The same trace with
work/output bit zero gives \(M_{=}(0,1)=0\).
The tape-by-tape rendering in Figure~\ref{fig:eq-machine-trace} keeps every
scanned location visible through the full run.

\begin{figure}[t]
\centering
\begin{tikzpicture}
  \matrix (trace) [matrix of nodes,
    nodes={cell, anchor=center},
    row sep=3.2mm,
    column sep=-\pgflinewidth,
    ampersand replacement=\&] {
    |[draw=none,fill=none,lblsm,minimum width=15mm]| time \&
    |[draw=none,fill=none,lblsm]| state \&
    |[draw=none,fill=none,tapelbl,anchor=center]| $A_0$ \&
    |[draw=none,fill=none,tapelbl,anchor=center]| $A_1$ \&
    |[draw=none,fill=none,minimum width=2.5mm]| \&
    |[draw=none,fill=none,tapelbl,anchor=center]| $B_0$ \&
    |[draw=none,fill=none,tapelbl,anchor=center]| $B_1$ \&
    |[draw=none,fill=none,minimum width=2.5mm]| \&
    |[draw=none,fill=none,tapelbl,anchor=center]| $W_0$ \&
    |[draw=none,fill=none,tapelbl,anchor=center]| $W_1$ \&
    |[draw=none,fill=none,minimum width=2.5mm]| \&
    |[draw=none,fill=none,tapelbl,anchor=center]| $O_0$ \&
    |[draw=none,fill=none,tapelbl,anchor=center]| $O_1$ \&
    |[draw=none,fill=none,lblsm,minimum width=28mm,anchor=west]| rule used \\
    |[draw=none,fill=none,lblsm,minimum width=15mm]| $t=0$ \&
    |[state initial]| $q_0$ \&
    |[cell head]| $1$ \& |[cell blank]| $\sqcup$ \&
    |[draw=none,fill=none,minimum width=2.5mm]| \&
    |[cell head]| $1$ \& |[cell blank]| $\sqcup$ \&
    |[draw=none,fill=none,minimum width=2.5mm]| \&
    |[cell head]| $\sqcup$ \& |[cell blank]| $\sqcup$ \&
    |[draw=none,fill=none,minimum width=2.5mm]| \&
    |[cell head]| $\sqcup$ \& |[cell blank]| $\sqcup$ \&
    |[draw=none,fill=none,lblsm,minimum width=28mm,anchor=west]| equal-input row \\
    |[draw=none,fill=none,lblsm,minimum width=15mm]| $t=1$ \&
    |[state]| $q_1$ \&
    |[cell]| $1$ \& |[cell head]| $\sqcup$ \&
    |[draw=none,fill=none,minimum width=2.5mm]| \&
    |[cell]| $1$ \& |[cell head]| $\sqcup$ \&
    |[draw=none,fill=none,minimum width=2.5mm]| \&
    |[cell head]| $1$ \& |[cell blank]| $\sqcup$ \&
    |[draw=none,fill=none,minimum width=2.5mm]| \&
    |[cell head]| $\sqcup$ \& |[cell blank]| $\sqcup$ \&
    |[draw=none,fill=none,lblsm,minimum width=28mm,anchor=west]| copy row, $r=1$ \\
    |[draw=none,fill=none,lblsm,minimum width=15mm]| $t=2$ \&
    |[state initial]| $q_0$ \&
    |[cell]| $1$ \& |[cell head]| $\sqcup$ \&
    |[draw=none,fill=none,minimum width=2.5mm]| \&
    |[cell]| $1$ \& |[cell head]| $\sqcup$ \&
    |[draw=none,fill=none,minimum width=2.5mm]| \&
    |[cell head]| $1$ \& |[cell blank]| $\sqcup$ \&
    |[draw=none,fill=none,minimum width=2.5mm]| \&
    |[cell head]| $1$ \& |[cell blank]| $\sqcup$ \&
    |[draw=none,fill=none,lblsm,minimum width=28mm,anchor=west]| halt row, $r=1$ \\
    |[draw=none,fill=none,lblsm,minimum width=15mm]| $t=3$ \&
    |[state halt]| $q_{\mathrm h}$ \&
    |[cell]| $1$ \& |[cell head]| $\sqcup$ \&
    |[draw=none,fill=none,minimum width=2.5mm]| \&
    |[cell]| $1$ \& |[cell head]| $\sqcup$ \&
    |[draw=none,fill=none,minimum width=2.5mm]| \&
    |[cell head]| $1$ \& |[cell blank]| $\sqcup$ \&
    |[draw=none,fill=none,minimum width=2.5mm]| \&
    |[cell head]| $1$ \& |[cell blank]| $\sqcup$ \&
    |[draw=none,fill=none,lblsm,minimum width=28mm,anchor=west]| halted \\
  };
  \foreach \r/\c in {2/3,2/6,2/9,2/12,3/4,3/7,3/9,3/12,
                     4/4,4/7,4/9,4/12,5/4,5/7,5/9,5/12}
    \node[head, below=.5mm of trace-\r-\c] {};
  \draw[brace] ([xshift=-3mm]trace-2-1.north west) --
    ([xshift=-3mm]trace-5-1.south west)
    node[midway, left=4mm, lblsm, align=right] {three\\transitions};
  \node[checked, below=5mm of trace-5-9, minimum width=47mm]
    {$M_{=}(1,1)=1\qquad \Time_{M_=,(1,1)}=3$};
\end{tikzpicture}

\caption{Four time rows track the \(A,B,W,O\) tapes of \(M_=\), with amber cells and triangles marking every scanned symbol and head.  The rule labels expose why equal inputs produce output one and reach \(q_{\mathrm h}\) after exactly three transitions.}
\label{fig:eq-machine-trace}
\end{figure}

\paragraph{Why this matters here.}
The example is intentionally mundane: a paper verifier is just a much larger
finite transition table whose output bit says accept or reject.  Statements
about a verifier's code, time, or simulation refer to this concrete object,
even when the paper describes it in higher-level pseudocode.

\subsection{Several inputs and the verifier convention}

Writing \(M(x_1,\ldots,x_k)\) means that \(x_i\) initially occupies the
initial cells of input tape \(i\), with all other cells blank and all heads at
zero.  Multiple tapes are convenient interfaces, not a stronger model: a
one-tape machine can store a self-delimiting encoding of all tapes and their
head markers and simulate one multitape step with polynomial overhead.

The paper's \emph{decider} is a total five-input machine
\[
 D(n,x,y,a,b)\in\{0,1\}.
\]
Here \(n\) is an index written in binary, \(x,y\) are the two questions, and
\(a,b\) are the two answers.  Total means that \(D\) halts on every such
five-tuple.  Its paper-level time at index \(n\) is the worst case over the
other four strings.  A normal-form verifier is a pair \(V=(S,D)\), where the
sampler \(S\) produces the questions and \(D\) checks the answers.
Figure~\ref{fig:verifier-normal-form} draws that interface without bundling
away any of the decider's five inputs.

\begin{figure}[t]
\centering
\begin{tikzpicture}[node distance=8mm and 17mm]
  \node[ctrl, minimum width=23mm] (sampler) {sampler $S$};
  \node[stage, above right=7mm and 28mm of sampler, minimum width=22mm] (pa)
    {prover A};
  \node[stage, below right=7mm and 28mm of sampler, minimum width=22mm] (pb)
    {prover B};
  \node[ctrl, right=76mm of sampler, minimum width=35mm, minimum height=18mm] (decider)
    {decider $D$\\[-1pt]$D(n,x,y,a,b)\in\{0,1\}$};

  \coordinate (xfork) at ($(sampler.east)!0.30!(pa.west)$);
  \coordinate (yfork) at ($(sampler.east)!0.30!(pb.west)$);
  \draw[->, machine, line width=.8pt] (sampler.east) |- node[lblsm, above, pos=.70] {$x$} (pa.west);
  \draw[->, machine, line width=.8pt] (sampler.east) |- node[lblsm, below, pos=.70] {$y$} (pb.west);
  \draw[->, machine, line width=.8pt] (pa.east) -| node[lblsm, above, pos=.25] {$a$}
    ([yshift=4mm]decider.west);
  \draw[->, machine, line width=.8pt] (pb.east) -| node[lblsm, below, pos=.25] {$b$}
    ([yshift=-4mm]decider.west);
  \draw[machine, line width=.7pt] (xfork) |- ([yshift=7mm]decider.west)
    node[lblsm, above, pos=.82] {$x$};
  \draw[machine, line width=.7pt] (yfork) |- ([yshift=-7mm]decider.west)
    node[lblsm, below, pos=.82] {$y$};
  \fill[machine] (xfork) circle (1.2pt);
  \fill[machine] (yfork) circle (1.2pt);
  \draw[->, machine, line width=.8pt] ([yshift=15mm]decider.north) --
    node[lblsm, right] {$n$ in binary} (decider.north);

  \node[cell head, right=11mm of decider] (bit) {$0/1$};
  \draw[->, machine, line width=.8pt] (decider) -- node[lblsm, above] {accept bit} (bit);
  \node[checked, below=9mm of decider, minimum width=36mm] (total)
    {halts on every five-tuple};
  \draw[callout] (total) -- (decider);
  \begin{scope}[on background layer]
    \node[panel, fit=(sampler)(pa)(pb)(decider)(total), inner sep=6mm] (verifier) {};
  \end{scope}
  \node[lbl, text=machine, anchor=south west] at ([xshift=2mm,yshift=1mm]verifier.north west)
    {$V=(S,D)$};
\end{tikzpicture}

\caption{The sampler sends \(x,y\) to the two provers, while the total decider receives the copied questions together with \(n\) and returned answers \(a,b\).  This five-port interface is the normal form preserved by the later source transformations.}
\label{fig:verifier-normal-form}
\end{figure}

\paragraph{Why this matters here.}
The five inputs are not later bundled by sleight of hand.  When we translate a
decider to a lambda term, its argument is a five-component encoded tuple; when
we translate back, those components initialize five separate tapes.  The
worst-case convention is what lets a resource bound hold independently of
unread suffixes of arbitrarily long answer strings.

\paragraph{Worked verifier example.}
Let \(S\) always ask the one-bit questions \(x=y=0\), and let \(D\) ignore
\(n,x,y\) and run \(M_=\) on the first bits of \(a,b\).  Then the players are
accepted exactly when their one-bit answers agree, and the decision stage
takes three transitions plus fixed input-parsing overhead.

\subsection{Descriptions are finite data}

Because \(Q,\Gamma\), and \(\delta\) are finite, we may serialize them with a
fixed prefix-free convention.  The resulting bit string is the machine
description \(\langle M\rangle\); the paper also writes \(\overline M\).
Its length in bits is
\[
 |M|:=|\langle M\rangle|.
\]
Malformed strings are assigned a fixed rejecting machine, so every bit string
still has a defined interpretation.  The encoding is syntactic: extensionally
equivalent machines may have different codes, sizes, and running times.  A
useful physics analogy is that a description is like a written Hamiltonian,
not its spectrum.  The actual definition is the serialized state set and
transition table, and an algorithm may inspect those bits without solving the
machine's behavior.
Figure~\ref{fig:description-bits} makes the syntax--behavior distinction
concrete for the five-row equality machine.

\begin{figure}[t]
\centering
\begin{tikzpicture}[node distance=8mm and 0mm]
  \node[descbox, minimum width=20mm, minimum height=13mm] (machine)
    {$M_=$\\[-1pt]\scriptsize three states, five rows};
  \node[bytes, right=13mm of machine, minimum width=14mm] (qbits) {001\,11};
  \node[bytes, right=0pt of qbits, minimum width=14mm] (gbits) {010\,01};
  \node[bytes, right=0pt of gbits, minimum width=15mm] (d1) {00\,1101};
  \node[bytes, right=0pt of d1, minimum width=15mm] (d2) {01\,0011};
  \node[bytes, right=0pt of d2, minimum width=15mm] (d3) {10\,1010};
  \node[bytes, right=0pt of d3, minimum width=15mm] (d4) {11\,0110};
  \node[bytes, right=0pt of d4, minimum width=15mm] (d5) {00\,0001};
  \draw[quote] (machine) -- node[lblsm, above, text=desc] {serialize} (qbits);
  \node[byteslbl, below=2mm of qbits] {$Q$};
  \node[byteslbl, below=2mm of gbits] {$\Gamma$};
  \node[byteslbl, below=2mm of d1] {$\delta_1$};
  \node[byteslbl, below=2mm of d2] {$\delta_2$};
  \node[byteslbl, below=2mm of d3] {$\delta_3$};
  \node[byteslbl, below=2mm of d4] {$\delta_4$};
  \node[byteslbl, below=2mm of d5] {$\delta_5$};
  \draw[brace] ([yshift=9mm]qbits.north west) -- ([yshift=9mm]d5.north east)
    node[midway, above=4mm, lbl, text=desc] {prefix-free $\langle M_=\rangle$};

  \node[descbox, below=19mm of d2, xshift=-21mm, text width=43mm] (ham)
    {written Hamiltonian\\[-1pt]\scriptsize finite coefficients};
  \node[bad, rounded corners=2pt, right=18mm of ham, text width=43mm,
        align=center, inner sep=5pt] (spectrum)
    {spectrum\\[-1pt]\scriptsize full behavior};
  % thin callouts, started below the field labels so they do not touch them;
  % the crossed dashed one says the spectrum side is simply not there.
  \draw[callout] ([yshift=-6mm]d2.south) -- (ham.north);
  \draw[callout, bad, dash pattern=on 2pt off 2pt]
    ([yshift=-6mm]d4.south) -- (spectrum.north);
  \node[text=bad, fill=white, inner sep=1pt, font=\footnotesize\bfseries]
    at ($([yshift=-6mm]d4.south)!0.5!(spectrum.north)$) {$\times$};
  \node[lblsm, below=3mm of ham] {syntax available to an algorithm};
  \node[lblsm, below=3mm of spectrum, text=bad] {not encoded extensionally};
\end{tikzpicture}

\caption{A prefix-free strip stores the state alphabet and all five transition rows of \(M_=\), with field boundaries remaining available to a parser.  Like a written Hamiltonian rather than its spectrum, \(\langle M_=\rangle\) records finite syntax rather than the machine's full behavior.}
\label{fig:description-bits}
\end{figure}

\paragraph{Why this matters here.}
The size bound \(|V|\leq\lambda\) constrains the verifier's source code, not
the number of functions extensionally equivalent to it.  Transformations such
as introspection and answer reduction copy, scan, hardwire, and simulate that
source.  They cannot be defined merely on a black-box input/output relation.

\paragraph{Worked description example.}
The code of \(M_=\) contains tags for the three states and the five rows above.
A compiler can prepend a wrapper saying ``run this code, but reject after
three simulated steps.''  It produces new finite code without having to
decide any semantic property of \(M_=\).
The compilation itself is drawn in Figure~\ref{fig:timeout-wrapper}.

\begin{figure}[t]
\centering
\begin{tikzpicture}[node distance=9mm and 12mm]
  \node[descbox, minimum width=25mm] (old) {$\langle M_=\rangle$};
  \node[descbox, below=of old, minimum width=25mm, fill=focuslt, draw=focus] (bound)
    {$T=3$};
  \node[checked, right=18mm of old, yshift=-8mm, minimum width=26mm] (compiler)
    {wrapper\\compiler};
  \draw[xform] (old) -- (compiler);
  \draw[xform] (bound) -- (compiler);

  \node[descbox, right=18mm of compiler, minimum width=66mm, minimum height=35mm] (wrapped) {};
  \draw[xform] (compiler) -- node[xformlbl] {emits} (wrapped.west);
  \node[lbl, text=desc, anchor=north west] at ([xshift=3mm,yshift=-3mm]wrapped.north west)
    {new finite description};
  \node[descbox, minimum width=25mm] (inner) at ([xshift=-14mm,yshift=-3mm]wrapped.center)
    {$\langle M_=\rangle$};
  \node[checked, right=5mm of inner, text width=25mm] (clock)
    {count steps\\$0,1,2,3$};
  \draw[->, machine, line width=.8pt] (inner) -- (clock);
  \node[bad, rounded corners=1.5pt, below=3mm of clock, inner sep=3pt,
        font=\scriptsize\sffamily] (reject) {reject if still running};
  \draw[->, bad, line width=.8pt] (clock) -- (reject);

  \node[lbl, below=7mm of wrapped, align=center]
    {copies source and a literal counter\\without asking whether \(M_=\) would halt};
\end{tikzpicture}

\caption{The compiler copies \(\langle M_=\rangle\) beside a literal three-step counter whose exhausted branch rejects.  Producing this bounded wrapper is a terminating syntactic operation and never requires deciding the wrapped machine's unbounded behavior.}
\label{fig:timeout-wrapper}
\end{figure}

A verifier description is the ordered pair
\(\langle V\rangle=(\langle S\rangle,\langle D\rangle)\).  In
\code{fig:compress}, \(\Compress\) literally receives
\((\langle V\rangle,\lambda)\), constructs descriptions for introspection,
answer reduction, and parallel repetition, and returns the description of a
new verifier.  This is why it is meaningful for \(\Compress\) to take ``a
verifier'' as input: it takes the finite string naming the verifier's rules,
not an infinite table of all possible plays.

\section{Universality: one machine can interpret all machines}
\label{sec:refresher-universality}

\subsection{The universal machine}

Fix an encoding of tuples and programs.  A universal machine \(U_k\) is a
single machine satisfying
\[
 U_k(\langle M\rangle,x_1,\ldots,x_k)=M(x_1,\ldots,x_k)
\]
whenever the right side halts, and diverging when that simulated run diverges.
Operationally, \(U_k\) keeps the code of \(M\) on a read-only region, stores an
encoded simulated configuration, finds the matching transition-table entry,
and updates the encoding.  Universality therefore says that descriptions can
be used as ordinary input data.

The paper needs the quantitative version.  With its dual-rail tuple encoding,
there are universal constants \(C_U,c_U\geq1\) such that a \(T\)-step run of a
\(k\)-input machine is reproduced within
\[
 C_U\bigl(k|\langle M\rangle||x|T\bigr)^{c_U},
 \qquad |x|=\sum_i|x_i|.
\]
No particular exponent is sacred.  What matters is that the overhead is one
fixed polynomial after the encodings have been fixed.
Figure~\ref{fig:universal-machine} exposes the read--lookup--update loop behind
this quantitative simulation statement.

\begin{figure}[t]
\centering
\begin{tikzpicture}[node distance=5mm and 8mm]
  \node[bytes, minimum width=17mm] (code0) {01101};
  \node[bytes, right=0pt of code0, minimum width=17mm] (code1) {11000};
  \node[bytes, right=0pt of code1, minimum width=17mm] (code2) {10110};
  \node[bytes, right=0pt of code2, minimum width=17mm] (code3) {00101};
  \node[byteslbl, above=2mm of code0.north west, anchor=south west] (codelbl)
    {read-only program region \(\langle M\rangle\)};
  \node[cell, right=14mm of code3] (cq) {$q$};
  \node[cell, right=0pt of cq] (cs0) {$0$};
  \node[cell head, right=0pt of cs0] (cs1) {$1$};
  \node[cell blank, right=0pt of cs1] (cs2) {$\sqcup$};
  \node[cell, right=0pt of cs2] (ch) {$h$};
  \node[head, below=.8mm of cs1] {};
  \node[tapelbl, above=2mm of cq.north west, anchor=south west]
    {encoded simulated configuration};

  \node[ctrl, below=18mm of code2, xshift=18mm, minimum width=31mm] (uk)
    {universal control\\$U_k$};
  \node[checked, left=10mm of uk, text width=27mm] (lookup)
    {1.\ scan code\\2.\ match \(\delta\)-row};
  \node[checked, right=10mm of uk, text width=27mm] (update)
    {3.\ rewrite config\\4.\ advance time};
  \draw[->, desc, line width=.9pt] (code2.south) -- (lookup.north);
  \draw[->, machine, line width=.9pt] (update.north) -- (cs2.south);
  \draw[xform] (lookup) -- (uk);
  \draw[xform] (uk) -- (update);
  \draw[->, focus, line width=1.1pt, rounded corners=5pt]
    (update.south) -- ++(0,-7mm) -| node[lblsm, below, pos=.27, text=focus]
    {repeat once per simulated step} (lookup.south);

  \node[stage, below=17mm of uk, minimum width=100mm] (cost)
    {$T\ \text{steps of }M
      \quad\Longrightarrow\quad
      C_U\bigl(k|\langle M\rangle|\,|x|\,T\bigr)^{c_U}\ \text{steps of }U_k$};
  \node[lblsm, right=4mm of cost] {$|x|=\sum_i|x_i|$};
\end{tikzpicture}

\caption{The fixed control of \(U_k\) alternates between a read-only program region and the encoded simulated configuration, matching one \(\delta\)-row before each update.  Iterating this real lookup loop yields the single fixed-polynomial overhead used by generated verifiers.}
\label{fig:universal-machine}
\end{figure}

\paragraph{Why this matters here.}
Generated verifiers call old verifiers as subroutines.  Their finite source
contains the universal interpreter plus the old description; it does not
contain a separate hand-compiled copy of every possible transition.  The
polynomial overhead is later absorbed into a new resource exponent.

\paragraph{Worked universal simulation.}
On input \((\langle M_=\rangle,1,0)\), \(U_2\) parses the five table rows,
builds the initial four-tape configuration, and simulates the three transitions
above.  It outputs zero.  The interpreter is fixed; changing the two input
bits or replacing \(\langle M_=\rangle\) changes only its data.

\subsection{Parameter fixing: the
\texorpdfstring{\(s\)-\(m\)-\(n\)}{s-m-n} theorem}

The notation \(s\)-\(m\)-\(n\) packages a simple compiler operation.  If
program \(e\) expects \(r+k\) arguments, there is a total computable function
\(s_{r,k}\) on bit strings such that
\[
 \varphi^{(k)}_{s_{r,k}(e,z_1,\ldots,z_r)}(x_1,\ldots,x_k)
 =\varphi^{(r+k)}_e(z_1,\ldots,z_r,x_1,\ldots,x_k).
\]
The left program has the first \(r\) values hardwired into its source.

\paragraph{Proof idea in two lines.}
Print a wrapper containing one literal copy of \(e,z_1,\ldots,z_r\).  On input
\(x\), the wrapper forms the tuple
\((e,z_1,\ldots,z_r,x)\) and invokes the universal machine; printing this
wrapper is itself a terminating, effective source transformation.
Figure~\ref{fig:smn} shows exactly which slots become source literals and
which remain live ports.

\begin{figure}[t]
\centering
\begin{tikzpicture}[node distance=7mm and 8mm]
  \node[lbl, text=desc, font=\footnotesize\sffamily\bfseries] (source)
    {SOURCE WITH \(r+k\) SLOTS};
  \node[descbox, below=3mm of source, minimum width=52mm] (program)
    {$e(\ \underbrace{z_1,\ldots,z_r}_{\text{static}}\ ;\
          \underbrace{x_1,\ldots,x_k}_{\text{live}}\ )$};
  \node[descbox, below=5mm of program, fill=focuslt, draw=focus] (lits)
    {literal data \(z_1,\ldots,z_r\)};
  \node[checked, right=10mm of program, yshift=-7mm, minimum width=26mm] (smn)
    {$s_{r,k}$\\compiler};
  \draw[xform] (program) -- (smn);
  \draw[xform] (lits) -- (smn);

  % 50mm: the 42mm literal strip inset by 4mm must not spill over the panel
  \node[descbox, right=10mm of smn, minimum width=50mm, minimum height=31mm] (wrapper) {};
  \draw[xform] (smn) -- (wrapper.west);
  \node[bytes, anchor=north west, minimum width=42mm]
    at ([xshift=4mm,yshift=-5mm]wrapper.north west) (literal)
    {\(e\mid z_1\mid\cdots\mid z_r\)};
  \node[ctrl, below=4mm of literal, minimum width=42mm] (univ)
    {invoke \(U_{r+k}\)};
  \node[stage, below=9mm of wrapper, minimum width=51mm] (live)
    {run-time inputs \(x_1,\ldots,x_k\)};
  \draw[->, machine, line width=.8pt] (live) -- (univ);
  \draw[->, desc, line width=.8pt] (literal) -- (univ);
  \node[byteslbl, above=2mm of wrapper] {wrapper code \(s_{r,k}(e,z_1,\ldots,z_r)\)};

  \node[checked, below=10mm of live, xshift=-31mm, minimum width=109mm,
        font=\scriptsize] (law)
    {$\varphi^{(k)}_{s_{r,k}(e,z_1,\ldots,z_r)}(x_1,\ldots,x_k)$\\[-1pt]
     $=\varphi^{(r+k)}_e(z_1,\ldots,z_r,x_1,\ldots,x_k)$};
\end{tikzpicture}

\caption{The \(s_{r,k}\) compiler moves \(e,z_1,\ldots,z_r\) into a literal byte region while preserving \(x_1,\ldots,x_k\) as run-time inputs.  The resulting wrapper realizes parameter fixing by one finite source transformation followed by universal simulation.}
\label{fig:smn}
\end{figure}

\paragraph{Why this matters here.}
Hardwiring is the bridge from ``this value is supplied at run time'' to ``this
value is part of the verifier description.''  Compression constantly fixes a
verifier, a level, or a time exponent while leaving \((n,x,y,a,b)\) as the
live inputs.  Part~II calls the corresponding operation \(\Specialize\).

\paragraph{Worked specialization.}
Apply \(s_{1,1}\) to the two-input equality machine and the static bit \(1\).
The resulting one-input program has code
\(e_1=s_{1,1}(\langle M_=\rangle,1)\) and satisfies
\(\varphi_{e_1}(a)=M_=(1,a)\).  Thus it is the one-bit identity test, even
though its source is a universal-simulation wrapper rather than a newly
invented transition table.

\subsection{Why the Halting problem is undecidable}

Let
\[
 \mathsf{HALT}=\{(\langle M\rangle,x):M(x)\neq\bot\}.
\]
The Halting problem asks for a total decider \(H\) for this set.  The obstruction
is diagonalization, and it is worth seeing every step.

\begin{proof}[Diagonal proof]
Assume, for contradiction, that \(H(e,x)\) always halts and returns one exactly
when the program described by \(e\) halts on \(x\).  Construct a one-input
machine \(D\) as follows on input \(z\):
\begin{enumerate}
\item Run \(H(z,z)\).
\item If \(H\) returns zero, halt.
\item If \(H\) returns one, loop forever.
\end{enumerate}
This is a legitimate finite program because \(H\) was assumed to be a total
subroutine.  Let \(d=\langle D\rangle\) and run \(D(d)\).

If \(H(d,d)=0\), then by the specification of \(H\), \(D(d)\) does not halt;
but Step 2 says that it halts.  If \(H(d,d)=1\), then the specification says
that \(D(d)\) halts; but Step 3 says that it loops forever.  The two possible
output bits both yield contradictions.  Therefore no such total \(H\) exists.
\end{proof}

Figure~\ref{fig:diagonal} places the two contradictory cases directly beside
the self-input entry that \(D\) flips.

\begin{figure}[t]
\centering
\begin{tikzpicture}[node distance=7mm and 12mm]
  \matrix (halt) [matrix of nodes,
    nodes={config, minimum width=18mm, minimum height=8mm, anchor=center},
    row sep=-\pgflinewidth, column sep=-\pgflinewidth,
    ampersand replacement=\&] {
    |[draw=none,fill=none,lblsm]| $H(e,x)$ \&
    |[draw=none,fill=none,lblsm]| $x=e_1$ \&
    |[draw=none,fill=none,lblsm]| $x=e_2$ \&
    |[draw=none,fill=none,lblsm]| $x=d$ \\
    |[draw=none,fill=none,lblsm]| $e=e_1$ \&
    |[fill=focuslt,draw=focus]| $H(e_1,e_1)$ \&
    $H(e_1,e_2)$ \& $H(e_1,d)$ \\
    |[draw=none,fill=none,lblsm]| $e=e_2$ \&
    $H(e_2,e_1)$ \&
    |[fill=focuslt,draw=focus]| $H(e_2,e_2)$ \&
    $H(e_2,d)$ \\
    |[draw=none,fill=none,lblsm]| $e=d$ \&
    $H(d,e_1)$ \& $H(d,e_2)$ \&
    |[fill=focuslt,draw=focus,line width=1pt]| $H(d,d)$ \\
  };
  \draw[brace] ([yshift=4mm]halt-2-2.north west) --
    ([yshift=4mm]halt-4-4.north east)
    node[midway, above=4mm, lbl, text=focus] {self-input diagonal};

  \node[checked, right=16mm of halt-4-4, text width=31mm] (flip)
    {$D(z)$ asks $H(z,z)$\\[2pt]
     \(0\mapsto\) halt\\
     \(1\mapsto\) loop};
  \draw[->, focus, line width=1.2pt] (halt-4-4) -- node[lblsm, above, text=focus]
    {flip} (flip);
  \node[bad, rounded corners=2pt, below left=9mm and 7mm of flip,
        text width=29mm, align=center, inner sep=4pt] (zero)
    {$H(d,d)=0$\\says \emph{does not halt}\\but \(D(d)\) halts};
  \node[bad, rounded corners=2pt, below right=9mm and 7mm of flip,
        text width=29mm, align=center, inner sep=4pt] (one)
    {$H(d,d)=1$\\says \emph{halts}\\but \(D(d)\) loops};
  \draw[->, bad, line width=.8pt] (flip) -- (zero);
  \draw[->, bad, line width=.8pt] (flip) -- (one);
\end{tikzpicture}

\caption{The amber diagonal specializes \(H(e,x)\) to self-inputs, and \(D\) reverses the highlighted value \(H(d,d)\).  Both possible bits terminate in a red contradiction, so no total table-filling decider \(H\) can exist.}
\label{fig:diagonal}
\end{figure}

This proof does not say that a particular long run can never be understood.
It says there is no single terminating algorithm which answers the halting
question correctly for every program-description/input pair.

\paragraph{Why this matters here.}
The \(\mathrm{MIP}^*=\mathrm{RE}\) construction reduces the behavior of an
arbitrary machine to the value of a nonlocal game.  Its Halting verifier is
given \(\langle M\rangle\), embeds and simulates that description, and asks
whether bounded prefixes of the computation are consistent.  It does not
contain an impossible halting decider; rather, the infinite family of game
instances reflects the recursively enumerable process of eventually seeing a
halt.

\paragraph{Worked semidecision.}
For a fixed \(M,x\), simulate one more step at each stage \(t=0,1,2,\ldots\).
If \(M(x)\) halts after \(T\) steps, stage \(T\) witnesses it.  If it never
halts, this search itself never returns.  That one-sided behavior is the
``RE'' side of the theorem's computability interface.

Figure~\ref{fig:semidecision} separates the finite yes-witness from the
deliberately absent no-return.

\begin{figure}[t]
\centering
\begin{tikzpicture}[node distance=6mm and 8mm]
  \node[lbl, text=machine, font=\footnotesize\sffamily\bfseries] (yeslabel)
    {IF \(M(x)\) HALTS};
  \node[config, below=3mm of yeslabel] (y0) {$t=0$};
  \node[config, right=of y0] (y1) {$t=1$};
  \node[config, right=of y1] (y2) {$t=2$};
  \node[lbl, right=of y2] (ydots) {$\cdots$};
  \node[config, right=of ydots] (ytm) {$t=T-1$};
  \node[checked, right=of ytm, fill=focuslt, draw=focus, minimum width=29mm] (witness)
    {$t=T$\\halt witnessed};
  \foreach \a/\b in {y0/y1,y1/y2,y2/ydots,ydots/ytm,ytm/witness}
    \draw[->, machine, line width=.8pt] (\a) -- (\b);

  % second timeline left-aligned with the first
  \node[config, below=21mm of y0.south west, anchor=north west] (n0) {$t=0$};
  \node[lbl, text=machine, font=\footnotesize\sffamily\bfseries,
        above=2mm of n0.north west, anchor=south west] (nolabel)
    {IF \(M(x)\) NEVER HALTS};
  \node[config, right=of n0] (n1) {$t=1$};
  \node[config, right=of n1] (n2) {$t=2$};
  \node[lbl, right=of n2] (ndots) {$\cdots$};
  \node[config, right=of ndots] (nt) {$t$};
  \node[config, right=of nt] (ntp) {$t+1$};
  \node[lbl, right=of ntp, text=machine] (more) {$\cdots$};
  \foreach \a/\b in {n0/n1,n1/n2,n2/ndots,ndots/nt,nt/ntp,ntp/more}
    \draw[->, machine, line width=.8pt] (\a) -- (\b);
  \draw[->, machine, line width=.8pt, opacity=.5] (more) -- ++(9mm,0);
  \draw[->, machine, line width=.8pt, opacity=.22]
    ([xshift=11mm]more.east) -- ++(9mm,0);
  \node[bad, rounded corners=2pt, below=5mm of ntp, text width=32mm,
        align=center, inner sep=4pt, font=\scriptsize\sffamily] (return)
    {the search has no negative return};
  \draw[callout] (return) -- (ntp);
  \node[descbox, below=16mm of ndots, minimum width=85mm] (re)
    {recursively enumerable: a finite accepting witness exists exactly on the yes side};
\end{tikzpicture}

\caption{Stage \(T\) supplies a finite halting witness on the upper timeline, while the lower search extends forever without producing a negative answer.  This asymmetric stopping behavior is exactly the recursively enumerable interface used by the Halting verifier.}
\label{fig:semidecision}
\end{figure}

\section{Programs which obtain their own descriptions}
\label{sec:refresher-kleene}

The diagonal argument used a program on its own code.  Kleene's second
recursion theorem turns that move into a general construction.  In its most
useful form here:

\begin{quote}
Every total computable transformation of programs has a program that behaves
like its own transform.
\end{quote}

More formally, if \(Q\) is a total computable map from descriptions of
\(k\)-input programs to descriptions of \(k\)-input programs, then there is a
description \(e_Q\) such that
\[
 \varphi^{(k)}_{e_Q}=\varphi^{(k)}_{Q(e_Q)}.
\]
This is behavioral equality, not necessarily literal equality of bit strings.

For the paper, this theorem buys a finite verifier description whose behavior
may depend on that very description.
\begin{theorem}[Kleene's second recursion theorem, refresher form]
\label{thm:refresher-kleene}
Every total computable transformation \(Q\) of \(k\)-input program codes has
an effectively constructible behavioral fixed point \(e_Q\).
\end{theorem}
The self-specialization square and its four equalities are laid out in
Figure~\ref{fig:kleene-square}.

\begin{figure}[t]
\centering
\begin{tikzpicture}[node distance=7mm and 9mm]
  \node[descbox, text width=48mm] (p)
    {$p(z,x)$ computes \(s(z,z)\), then \(Q(s(z,z))\),\\
     then universally simulates that code on \(x\)};
  \node[checked, right=14mm of p, minimum width=24mm] (s)
    {specialize\\\(s(p,p)\)};
  \node[descbox, right=14mm of s, minimum width=30mm,
        fill=focuslt, draw=focus] (eq)
    {$e_Q:=s(p,p)$};
  \draw[xform] (p) -- node[xformlbl] {two copies of \(p\)} (s);
  \draw[xform] (s) -- node[xformlbl] {emit code} (eq);
  \draw[->, desc, line width=.9pt, bend left=32] (eq.north) to
    node[lblsm, above, text=desc] {its own specialized description}
    (p.north);

  \node[stage, below=18mm of p, minimum width=26mm, font=\scriptsize] (r1)
    {$\varphi^{(k)}_{e_Q}(x)$};
  \node[stage, right=1mm of r1, minimum width=28mm, font=\scriptsize] (r2)
    {$\varphi^{(k)}_{s(p,p)}(x)$};
  \node[stage, right=1mm of r2, minimum width=28mm, font=\scriptsize] (r3)
    {$\varphi^{(k+1)}_p(p,x)$};
  \node[stage, right=1mm of r3, minimum width=33mm, font=\scriptsize] (r4)
    {$\varphi^{(k)}_{Q(s(p,p))}(x)$};
  \node[stage, right=1mm of r4, minimum width=29mm, font=\scriptsize] (r5)
    {$\varphi^{(k)}_{Q(e_Q)}(x)$};
  \foreach \a/\b in {r1/r2,r2/r3,r3/r4,r4/r5}
    \draw[<->, focus, line width=.9pt] (\a) -- node[lblsm, above, text=focus] {$=$} (\b);
  \draw[->, desc, dashed] (eq.south) -- (r1.north);
  \draw[->, machine, dashed] (p.south) -- (r3.north);
  \node[checked, below=7mm of r3, minimum width=64mm] (fixed)
    {$\varphi^{(k)}_{e_Q}=\varphi^{(k)}_{Q(e_Q)}$};
  \draw[callout] (fixed) -- (r3);
\end{tikzpicture}

\caption{Two literal copies of \(p\) enter specialization to create \(e_Q=s(p,p)\), whose run follows the five expressions joined by four amber equalities.  The commuting construction turns apparent circularity into finite code generation plus ordinary universal simulation.}
\label{fig:kleene-square}
\end{figure}

\begin{proof}
Use parameter fixing for one static and \(k\) live arguments.  Write \(s(z,w)\)
for the code obtained by fixing the first input of a suitable \((k+1)\)-input
program \(z\) to the literal \(w\); hence
\[
 \varphi^{(k)}_{s(z,w)}(x)=\varphi^{(k+1)}_z(w,x).
\]
Now define the partial computable function
\[
 g(z,x):=\varphi^{(k)}_{Q(s(z,z))}(x).
\]
There is a finite \((k+1)\)-input program \(p\) for \(g\): on \((z,x)\), it
computes \(s(z,z)\), runs the total transformer \(Q\) on that string, and
universally simulates the resulting program on \(x\).  Set
\(e_Q:=s(p,p)\).  Then, for every \(x\),
\[
\begin{aligned}
 \varphi^{(k)}_{e_Q}(x)
 &=\varphi^{(k)}_{s(p,p)}(x)\\
 &=\varphi^{(k+1)}_p(p,x)\\
 &=\varphi^{(k)}_{Q(s(p,p))}(x)\\
 &=\varphi^{(k)}_{Q(e_Q)}(x).
\end{aligned}
\]
If the last computation diverges, each displayed partial computation diverges
for the same reason; if it halts, all return the same output.  Thus the
partial functions are equal on every input.  Every source-manipulating step
used to construct \(e_Q\) terminates, so the fixed point is effective.
\end{proof}

\paragraph{Why this matters here.}
The Halting verifier must pass its own decider description into \(\Compress\).
The theorem says that this is ordinary program construction, not a circular
definition requiring an infinite source string.  It also separates two facts:
constructing the fixed-point code terminates, while running the partial
program denoted by that code may or may not terminate.

\paragraph{Worked quine-like example.}
Let \(Q(z)\) be the code of the zero-input program ``print the literal string
\(z\).''  The transformation \(Q\) is total because it merely prints a
wrapper.  Its fixed point \(e_Q\) satisfies
\(\varphi_{e_Q}=\varphi_{Q(e_Q)}\), so running \(e_Q\) prints \(e_Q\) itself.
The theorem has built a quine without requiring the source printer to know its
own final address in advance.
Figure~\ref{fig:quine} traces the identical finite string from code to output.

\begin{figure}[t]
\centering
\begin{tikzpicture}[node distance=8mm and 11mm]
  \node[descbox, minimum width=18mm, fill=focuslt, draw=focus] (z) {$z$};
  \node[checked, right=13mm of z, minimum width=24mm] (q)
    {$Q$\\source printer};
  \node[descbox, right=13mm of q, minimum width=43mm] (qz)
    {$Q(z)=\langle\texttt{print literal }z\rangle$};
  \draw[xform] (z) -- (q);
  \draw[xform] (q) -- (qz);
  \node[lbl, above=3mm of q] {compile time};

  \node[descbox, below=20mm of z, minimum width=24mm] (eq)
    {$\langle e_Q\rangle$};
  \node[ctrl, right=21mm of eq, minimum width=26mm] (run)
    {run \(e_Q\)};
  \node[cell head, right=16mm of run] (out0) {$0$};
  \node[cell, right=0pt of out0] (out1) {$1$};
  \node[cell, right=0pt of out1] (out2) {$1$};
  \node[cell, right=0pt of out2] (out3) {$0$};
  \node[cell blank, right=0pt of out3] (out4) {$\cdots$};
  \node[head, below=.8mm of out0] {};
  \draw[->, machine, line width=.8pt] (eq) -- (run);
  \draw[->, machine, line width=.8pt] (run) -- (out0);
  \draw[brace] ([yshift=4mm]out0.north west) -- ([yshift=4mm]out4.north east)
    node[midway, above=4mm, lblsm] {output tape \(=\langle e_Q\rangle\)};
  \draw[quote] (eq.west) to[bend left=52]
    node[lblsm, left, text=desc] {$z=e_Q=s(p,p)$} (z.west);
  % return arc kept ABOVE the tape: starting under it crossed the cells
  \draw[->, focus, line width=1.1pt, shorten >=1.5pt] (out0.north west)
    to[bend right=22]
    node[lblsm, above, text=focus] {same finite string} (eq.north east);
  \node[checked, below=10mm of run, minimum width=62mm]
    {$\varphi_{e_Q}=\varphi_{Q(e_Q)}$};
\end{tikzpicture}

\caption{Specializing the source printer at its own generated code produces \(e_Q\), whose output tape contains the same finite description \(\langle e_Q\rangle\).  The returning amber arrow emphasizes behavioral self-reproduction without requiring a source string of infinite regress.}
\label{fig:quine}
\end{figure}

\subsection{The paper's
\texorpdfstring{\(\mathcal F(\mathcal F,M,\lambda,\ldots)\)}
{F(F,M,lambda,...)} construction}

The paper writes the fixed point in an unrolled form.  Begin with an
eight-input machine
\[
 \mathcal F(f,m,\lambda,n,x,y,a,b).
\]
After its initial \(n\)-step simulation has not halted, its first source action
is computable: hardwire
\((f,m,\lambda)\) into the first three slots of \(\mathcal F\)'s template, obtaining
a five-input decider description \(d\).  It pairs \(d\) with a sampler
description, invokes \(\Compress\) on that verifier description, and runs the
resulting decider on \((n,x,y,a,b)\).  The published call is
\[
 \mathcal F(\langle\mathcal F\rangle,\langle M\rangle,\lambda,n,x,y,a,b).
\]
After the first three arguments have been fixed, the produced \(d\) is exactly
the description of the five-input program currently being run.  Hence the
verifier passed to \(\Compress\) contains its own decider code.
Figure~\ref{fig:halt-f-construction} follows both the bounded-halting shortcut
and the complete feedback path through \(\Compress\).

\begin{figure}[t]
\centering
\begin{tikzpicture}[node distance=6mm and 4mm]
  %% ---------------- top row: the eight-input machine and its early exit
  \node[descbox, minimum width=34mm] (inputs)
    {$\langle\mathcal F\rangle,\langle M\rangle,L$\\
     \(u=(n,x,y,a,b)\)};
  \node[ctrl, right=8mm of inputs, minimum width=24mm] (f)
    {$\mathcal F$\\eight inputs};
  \node[checked, right=8mm of f, minimum width=32mm] (precheck)
    {run \(M\) on blank\\for \(n\) steps};
  \node[checked, right=14mm of precheck, fill=focuslt, draw=focus,
        minimum width=21mm] (accept)
    {halt seen\\accept};
  \draw[xform] (inputs) -- (f);
  \draw[xform] (f) -- (precheck);
  \draw[->, check, line width=.9pt] (precheck) -- node[lblsm, above] {halts} (accept);

  %% ---------------- middle row: the executed source construction
  \node[checked, below=20mm of inputs, minimum width=24mm] (hardwire)
    {hardwire\\\((f,m,L)\)};
  \node[descbox, right=10mm of hardwire, minimum width=14mm,
        fill=focuslt, draw=focus] (d)
    {$\langle d\rangle$};
  \node[descbox, right=13mm of d, minimum width=25mm] (verifier)
    {$\langle V\rangle=$\\$(\langle S_L\rangle,\langle d\rangle)$};
  \node[checked, right=10mm of verifier, minimum width=25mm] (compress)
    {$\Compress$\\on \((\langle V\rangle,L)\)};
  \node[descbox, right=10mm of compress, minimum width=16mm] (dc)
    {$\langle d^{\mathrm c}\rangle$};
  \draw[xform] (hardwire) -- (d);
  \draw[->, focus, line width=1.2pt] (d) -- node[lblsm, above, text=focus]
    {self code} (verifier);
  \draw[xform] (verifier) -- (compress);
  \draw[xform] (compress) -- (dc);

  %% the "otherwise" branch runs in the band between the rows, not along the
  %% tops of the middle-row boxes
  \draw[->, check, line width=.9pt, rounded corners=3pt]
    (precheck.south) -- ++(0,-4mm)
    node[lblsm, right=0.8mm, pos=.55] {otherwise} -| (hardwire.north);
  % the produced self code IS the code of the machine now running
  \draw[->, focus, line width=1.1pt, shorten >=1pt] (d.north)
    to[out=95,in=-85]
    node[lblsm, right=2mm, pos=.5, text=focus, align=left, text width=40mm]
      {with \(f=\langle\mathcal F\rangle\), \(d\) is this running decider}
    (f.south);

  %% ---------------- what feeds the verifier, and what stays cited
  \node[checked, minimum width=30mm] (sampler)
    at ([yshift=-11mm]$(d.south)!0.5!(verifier.south)$)
    {$\ComputeSampler(L)$\\$\Downarrow\ \langle S_L\rangle$};
  \draw[->, desc, line width=.9pt] (sampler) -- (verifier);
  \node[citedbox, below=11mm of compress, minimum width=28mm] (theorem)
    {compression\\soundness and\\completeness};
  \draw[cited, dashed, ->] (theorem) -- (compress);

  %% ---------------- running the returned decider
  \node[checked, below=15mm of sampler, minimum width=38mm] (run)
    {run \(d^{\mathrm c}(n,x,y,a,b)\)};
  \node[cell head, below=6mm of run] (bit) {$0/1$};
  \draw[->, desc, line width=.9pt, rounded corners=3pt]
    (dc.south) |- (run.east);
  \draw[->, machine, line width=.8pt] (run) -- (bit);

  \node[tag check, anchor=north west, align=left]
    at ([yshift=-3mm]hardwire.south west)
    {EXECUTED\\SOURCE CONSTRUCTION};
  \node[tag cited, below=2mm of theorem]
    {CITED THEOREM};
\end{tikzpicture}

\caption{The eight-input machine first accepts a witnessed \(n\)-step halt; otherwise it hardwires \((f,m,\lambda)\), feeds the resulting self-code \(d\) into \((S_\lambda,d)\), compresses, and runs the returned decider.  Green boxes are executed source operations while the dashed slate box is the cited game-theoretic guarantee (after \code{fig:halt_f}).}
\label{fig:halt-f-construction}
\end{figure}

\paragraph{Why this matters here.}
This is Kleene's construction with the diagonal substitution exposed as
pseudocode.  The transformation \(Q\) says ``insert the candidate code into a
verifier, compress it, then run the returned decider.''  The
\(\mathcal F(\mathcal F,\ldots)\) trick supplies \(s(p,p)\) directly.  An
earlier version of the paper invoked Kleene's recursion theorem explicitly;
the published presentation changes the notation, not the computability fact.

\paragraph{Worked schematic expansion.}
Suppress the fixed parameters \(M,\lambda\) and denote the five live inputs by
\(u\).  For a proposed self-description \(f\), put
\(d_f=\operatorname{hardwire}(\langle\mathcal F\rangle,f)\).
Substituting \(f=\langle\mathcal F\rangle\) gives
\[
 d_{\langle\mathcal F\rangle}(u)
 =\mathcal F(\langle\mathcal F\rangle,u)
 =Q(d_{\langle\mathcal F\rangle})(u),
\]
which is the recursion-theorem equation with the compiler operation written
out rather than hidden under the symbol \(s\).
